% ================= B 卷 =================
% 难度定位：综合应用（跨章节组合，含易错点）
\juantou{B}
\defenBiaoFour{20}{32}{30}{18}

\dati{名词解释}{每小题 5 分，共 20 分}

\begin{ti}
\xiaoti{1.（5 分）联系}
\kongbai{2.6cm}
\end{ti}

\begin{ti}
\xiaoti{2.（5 分）否定之否定规律}
\kongbai{2.6cm}
\end{ti}

\begin{ti}
\xiaoti{3.（5 分）真理}
\kongbai{2.6cm}
\end{ti}

\begin{ti}
\xiaoti{4.（5 分）生产关系}
\kongbai{2.6cm}
\end{ti}

\dati{简答题}{每小题 8 分，共 32 分}

\begin{ti}
\xiaoti{1.（8 分）内容与形式是一对基本范畴，请说明二者相互制约的具体表现。}
\kongbai{4.4cm}
\end{ti}

\begin{ti}
\xiaoti{2.（8 分）感性认识与理性认识有何区别？二者又是怎样相互过渡的？}
\kongbai{4.4cm}
\end{ti}

\begin{ti}
\xiaoti{3.（8 分）生产力与生产关系的矛盾运动遵循什么规律？这一规律有何现实意义？}
\kongbai{4.4cm}
\end{ti}

\begin{ti}
\xiaoti{4.（8 分）为什么说上层建筑一定要适合经济基础的状况？不适合时会出现什么后果？}
\kongbai{4.4cm}
\end{ti}

\dati{论述题}{每小题 15 分，共 30 分}

\begin{ti}
\xiaoti{1.（15 分）试述矛盾普遍性与特殊性的辩证关系，并说明这一原理为什么是建设中国特色社会主义的哲学基础。}
\kongbai{9.5cm}
\end{ti}

\begin{ti}
\xiaoti{2.（15 分）实践是检验真理的唯一标准，但实践标准既是确定的又是不确定的。请论述这一论断的含义及其意义。}
\kongbai{9.5cm}
\end{ti}

\dati{材料分析题}{共 18 分，阅读所给材料后回答问题}

\begin{ti}
\xiaoti{1.（18 分）某科研团队承担一项关键材料攻关任务。最初按国外文献的工艺路线试制，连续失败；团队改为自行测试不同温度与配比下的性能，累计完成三百余组实验，逐步摸清规律，最终研制出性能达标的新材料。负责人说：“论文上的结论是别人的实践结果，能不能用，还得我们自己试。”\\

请回答：\\
（1）运用实践与认识的辩证运动规律，分析该团队的攻关过程。（9 分）\\
（2）说明科技创新中真理尺度与价值尺度的统一。（9 分）}
\kongbai{12.0cm}
\end{ti}
